\begin{table}[htbp]
    \centering
    \setcounter{table}{2}% mesmo número da tabela do roteiro
    \renewcommand{\arraystretch}{1.35}\setlength{\tabcolsep}{5pt}
    \begin{tabular}{l@{\hspace{1.6em}}rrr@{\hspace{1.6em}}r}
        \toprule
        métrica & exp. & teór. & \textcolor{black!55}{c/ carga} & desvio \\
        \midrule
        $\tau$ ($\mu$s) & 61 & 21 & \textcolor{black!55}{21} & 197,2\,\% \\
        $T_r$ ($\mu$s) & 66 & 46 & \textcolor{black!55}{46} & 43,9\,\% \\
        $T_s$ ($\mu$s) & 114 & 87 & \textcolor{black!55}{85} & 30,6\,\% \\
        $V_c$ (V) & 0,479 & 0,481 & \textcolor{black!55}{0,480} & $-$0,5\,\% \\
        \bottomrule
    \end{tabular}
    \par\smallskip{\footnotesize $T_r$ de 10\,\% a 90\,\%, $T_s$ na faixa de $\pm 2$\,\%, $\tau$ até 63,2\,\%; a saída medida foi suavizada por média móvel de 5,7\,$\mu$s.}
    \caption{RC com compensador avanço de fase ($C_{c}$ = 1 $\mu$F, $R_{c}$ = 1 k$\Omega$, $C_{1}$ = 47 nF, $C_{2}$ = 1 nF, $R_{1}$ = 15 k$\Omega$, $R_{2}$ = 390 k$\Omega$, $R$ = $R_6$ = 10 k$\Omega$).}
    \label{tab:tabela3}
\end{table}
